% !TEX root = ../../../main.tex
% 所属：第四部分：Fourier 分析 / 第2章：Fourier级数的收敛性
% 节：2.2 Fourier级数的逐点收敛
% 正文按原文件顺序保存；由 ../chapter2.tex 调用。
\section{Fourier级数的逐点收敛}

\begin{theorem}
	令 $f$ 是圆周上的可积函数,且在 $\theta_0$ 点可微.那么当 $N\to\infty$ 时,
	\[
	S_N(f)(\theta_0)\longrightarrow f(\theta_0).
	\]
\end{theorem}

\begin{proof}
	定义
	\[
	F(t)=
	\begin{cases}
		\displaystyle\frac{f(\theta_0-t)-f(\theta_0)}{t},
		& t\neq0,\ \abs{t}<\pi,\\[3mm]
		-f'(\theta_0),&t=0.
	\end{cases}
	\]
	首先,因为 $f$ 在 $\theta_0$ 点可微,所以 $F$ 在 $0$ 附近有界.其次,对充分小的 $\delta>0$,函数 $F$ 在
	\[
	[-\pi,-\delta]\cup[\delta,\pi]
	\]
	上可积；再由分段可积函数的拼接性质可知,$F$ 在整个区间 $[-\pi,\pi]$ 上可积.
	
	由
	\[
	S_N(f)(\theta_0)=(f\ast D_N)(\theta_0),
	\qquad
	\frac1{2\pi}\int_{-\pi}^{\pi}D_N(t)\dd t=1,
	\]
	可得
	\[
	\begin{aligned}
		S_N(f)(\theta_0)-f(\theta_0)
		&=\frac1{2\pi}\int_{-\pi}^{\pi}f(\theta_0-t)D_N(t)\dd t-f(\theta_0)\\
		&=\frac1{2\pi}\int_{-\pi}^{\pi}
		\bigl[f(\theta_0-t)-f(\theta_0)\bigr]D_N(t)\dd t\\
		&=\frac1{2\pi}\int_{-\pi}^{\pi}F(t)tD_N(t)\dd t.
	\end{aligned}
	\]
	再回顾 Dirichlet 核的公式
	\[
	tD_N(t)
	=\frac{t}{\sin(t/2)}
	\sin\bigl((N+\tfrac12)t\bigr),
	\]
	其中,$t/\sin(t/2)$ 在区间 $[-\pi,\pi]$ 上连续.又因为
	\[
	\sin\bigl((N+\dfrac12)t\bigr)
	=\sin(Nt)\cos(t/2)+\cos(Nt)\sin(t/2),
	\]
	所以
	\[
	\begin{aligned}
		S_N(f)(\theta_0)-f(\theta_0)
		={}&\frac1{2\pi}\int_{-\pi}^{\pi}
		F(t)\frac{t\cos(t/2)}{\sin(t/2)}\sin(Nt)\dd t\\
		&+\frac1{2\pi}\int_{-\pi}^{\pi}F(t)t\cos(Nt)\dd t.
	\end{aligned}
	\]
	上式中两个系数函数均为 Riemann 可积函数.对这两个积分应用 Riemann--Lebesgue 引理,便得
	\[
	S_N(f)(\theta_0)-f(\theta_0)\longrightarrow0.
	\]
\end{proof}

可以看出,上述证明中可微条件只用于保证函数
\[
F(t)=\frac{f(\theta_0-t)-f(\theta_0)}{t}
\]
在原点附近有界.定义
\[
F(t)=
\begin{cases}
	\dfrac{f(\theta_0-t)-f(\theta_0)}{t},
	& t\neq 0,\ |t|<\pi,\\[6pt]
	0,  & t=0.
\end{cases}
\]
因此,如果仅假设 $f$ 在 $\theta_0$ 点满足
Lipschitz 条件,定理的结论仍然成立.

事实上,设存在 $M>0$ 及 $\theta_0$ 的一个邻域,使得对该邻域内
任意的 $\theta$ 都有
\[
\abs{f(\theta)-f(\theta_0)}
\leq M\abs{\theta-\theta_0}.
\]
令 $\theta=\theta_0-t$,则当 $t$ 充分小时,
\[
\abs{F(t)}
=
\frac{\abs{f(\theta_0-t)-f(\theta_0)}}{\abs{t}}
\leq M.
\]
故 $F$ 在原点附近有界,其余证明与上述定理完全相同.
此时也称 $f$ 在 $\theta_0$ 点满足一阶 H\"older 条件.

\begin{corollary}
	假设 $f$ 和 $g$ 是定义在圆周上的两个可积函数,且对于某个 $\theta_0$,存在包含 $\theta_0$ 的开区间 $I$,满足
	\[
	f(\theta)=g(\theta),
	\qquad \theta\in I.
	\]
	那么当 $N\to\infty$ 时,
	\[
	S_N(f)(\theta_0)-S_N(g)(\theta_0)\longrightarrow0.
	\]
\end{corollary}

\begin{proof}
	令 $h=f-g$.因为 $h$ 在 $I$ 上恒为 $0$,所以 $h$ 在 $\theta_0$ 处可微,且 $h(\theta_0)=0$.由上一个定理,
	\[
	S_N(h)(\theta_0)\longrightarrow h(\theta_0)=0.
	\]
	再由 Fourier 部分和的线性性,
	\[
	S_N(h)(\theta_0)
	=S_N(f)(\theta_0)-S_N(g)(\theta_0),
	\]
	结论成立.
\end{proof}
